% Source: https://artofproblemsolving.com/community/c3h1602042_harmonic_mean
% Stable identity: aops-topic-1602042-q01
The configuration (drawn with $a>b$):

\begin{asy}
usepackage("amsmath");
size(7cm);
defaultpen(fontsize(10pt));
real a=6,bb=2.4; pair A=(0,0),B=(a+bb,0),O=((a+bb)/2,0),P=(a,0); real R=(a+bb)/2;
pair Q=(a,sqrt(a*bb)); pair H=O+dot(P-O,unit(Q-O))*unit(Q-O);
draw(arc(O,R,0,180),linewidth(0.9));draw(A--B,linewidth(0.9));
draw(O--Q,blue);draw(P--Q,deepgreen+linewidth(1));draw(P--H,red+dashed);
draw(H--Q,red+linewidth(1.6));
real q=0.25;draw(P+(-q,0)--P+(-q,q)--P+(0,q),deepgreen);
pair u=unit(O-H),v=unit(P-H);draw(H+q*u--H+q*(u+v)--H+q*v,red);
dot(A);dot(B);dot(O);dot(P);dot(Q);dot(H);
label("$A$",A,SW);label("$B$",B,SE);label("$O$",O,S);label("$P$",P,S);label("$Q$",Q,N);label("$H$",H,NW,red);
draw((0,-0.7)--(a,-0.7),Arrows(TeXHead));label("$a$",(a/2,-0.7),S);
draw((a,-0.7)--(a+bb,-0.7),Arrows(TeXHead));label("$b$",(a+bb/2,-0.7),S);
label("\mbox{}",(min(currentpicture,user=true)+max(currentpicture,user=true))/2,black); // keeps the text after the figure black when LaTeX typesets the labels inline
\end{asy}

\textbf{Arithmetic mean.} $OQ$ is a radius, and the diameter is $a+b$, so $OQ=\dfrac{a+b}{2}$.

\textbf{Geometric mean.} The angle $AQB$ is inscribed in a semicircle, so it is a right angle, and $QP$ is the altitude to the hypotenuse of the right triangle $AQB$. The triangles $APQ$ and $QPB$ are similar (each has a right angle at $P$, and $\angle PAQ=90^\circ-\angle PBQ=\angle PQB$), so $\dfrac{PQ}{a}=\dfrac{b}{PQ}$ and $PQ^2=ab$, i.e.\ $PQ=\sqrt{ab}$.

\textbf{Harmonic mean.} Triangle $OPQ$ has a right angle at $P$, and $PH$ is its altitude to the hypotenuse $OQ$. The right triangles $QHP$ and $QPO$ share the angle at $Q$, so they are similar and $\dfrac{HQ}{PQ}=\dfrac{PQ}{OQ}$. Hence
\[
HQ=\dfrac{PQ^2}{OQ}=\dfrac{ab}{\dfrac{a+b}{2}}=\dfrac{2ab}{a+b} .
\]

\textbf{The inequalities.} In the right triangle $OPQ$ the hypotenuse is the longest side, so $OQ\ge PQ$; in the right triangle $QHP$ (right angle at $H$) the hypotenuse $PQ$ is longer than the leg $HQ$. Hence
\[
\dfrac{a+b}{2}\ge\sqrt{ab}\ge\dfrac{2ab}{a+b},
\]
with strict inequalities when $a\ne b$ (then $O\ne P$, so $OQ>PQ$, and $H\ne Q$). $\blacksquare$

(Why the harmonic mean is natural: it is the reciprocal of the average of the reciprocals, $\dfrac{2ab}{a+b}=\dfrac{2}{\dfrac{1}{a}+\dfrac{1}{b}}$. It is the right ``average'' whenever reciprocals add, for example the average speed for two equal distances travelled at speeds $a$ and $b$. In the figure it appears because $HQ\cdot OQ=PQ^2$, i.e.\ $HQ\cdot\text{AM}=\text{GM}^2$.)
